\documentclass[11pt]{article}
\usepackage[a4paper,margin=1.9cm]{geometry}
\usepackage[T1]{fontenc}
\usepackage{textcomp}
\usepackage{amsmath,amssymb}
\usepackage{enumitem}
\usepackage{tikz}
\usetikzlibrary{calc,arrows.meta,patterns,decorations.pathmorphing}
\usepackage{xcolor}
\usepackage{multicol}
\usepackage{parskip}
\usepackage{lmodern}
\renewcommand{\familydefault}{\sfdefault}

\definecolor{unalblue}{RGB}{31,73,124}

\newlist{opciones}{enumerate}{1}
\setlist[opciones]{label=(\alph*),itemsep=1pt,topsep=2pt,leftmargin=1.8em,font=\normalfont}

\newcommand{\vect}[2]{\begin{pmatrix}#1\\#2\end{pmatrix}}
\newcommand{\vectiii}[3]{\begin{pmatrix}#1\\#2\\#3\end{pmatrix}}

\newcommand{\examheader}{%
  \noindent
  \begin{minipage}[t]{0.62\linewidth}
    {\bfseries\large Primer Parcial de C\'alculo en Varias Variables}\\
    {\small Septiembre de 2022}\\
    {\small Escuela de Matem\'aticas, Universidad Nacional de Colombia, Sede Medell\'in.}
  \end{minipage}%
  \hfill
  \begin{minipage}[t]{0.33\linewidth}
    \raggedleft\small
    Nombre: \rule{2.8cm}{0.4pt}\\[4pt]
    C\'edula: \rule{2.8cm}{0.4pt}
  \end{minipage}\\[4pt]
  \rule{\linewidth}{0.6pt}
}
\newcommand{\examfooter}{%
  \vfill
  \rule{\linewidth}{0.4pt}\\[1pt]
  {\scriptsize\textit{Transcripci\'on digital --- MathUNAL}\hfill\textcolor{unalblue}{\textbf{\thepage}}}
}
\pagestyle{empty}

\begin{document}

\examheader

\textbf{Instrucciones:} La duraci\'on del examen es de 1 hora y 45 minutos. Antes de empezar verifique que su examen est\'e completo y que sus datos sean correctos. No desengrape su examen ni a\~nada otras grapas. Use los espacios asignados para resolver cada pregunta. \textbf{No} est\'a permitido: hacer preguntas, usar calculadoras, sacar hojas en blanco, ni tomar apuntes en hojas diferentes a las del examen. Tampoco se permite el uso de celulares, ni smartphones, ni tabletas digitales.

\bigskip

\textit{En las preguntas 1 a 5 responda en el espacio en blanco. En estas preguntas s\'olo se tendr\'a en cuenta la respuesta al calificar.}

\bigskip
\noindent\textbf{1.} (24\,\%)

\begin{opciones}
  \item[(a)] (12\,\%) Sea $\displaystyle f(x,y) = \int_x^{xy} g(t)\,dt$, con $g:\mathbb{R}\to\mathbb{R}$ continua tal que $g(2)=-4$ y $g(8)=3$. Entonces
  $$f_x(2,4) = \underline{\hspace{3cm}}.$$
  \item[(b)] (12\,\%) La ecuaci\'on del plano tangente a la superficie dada por
  $$z = x\cos(x)\cos(y) + y$$
  en el punto $(0,\pi,\pi)$ es:
  $$\underline{\hspace{4cm}}.$$
\end{opciones}

\bigskip
\noindent\textbf{2.} (22\,\%)

\begin{opciones}
  \item[(a)] (10\,\%) La derivada direccional de la funci\'on $g(x,y,z)=z^2-x^2y$ en el punto $(1,1,2)$ en la direcci\'on del vector $\vec v = \vec\imath+2\vec\jmath+2\vec k$ es:
  $$\underline{\hspace{3cm}}.$$
  \item[(b)] (12\,\%) Las ecuaciones param\'etricas de la recta normal a la superficie $x^2+3y^2+2z^2=9$ en el punto $(2,1,1)$ son:
  $$x=\underline{\hspace{2.5cm}}, \qquad y=\underline{\hspace{2.5cm}}, \qquad z=\underline{\hspace{2.5cm}}.$$
\end{opciones}

\bigskip
\noindent\textbf{3.} (24\,\%)

\begin{opciones}
  \item[(a)] (12\,\%) Considere la siguiente ecuaci\'on:
  $$\cos(y+z) = 2x^2+e^{-z}+y^3.$$
  Al derivar la ecuaci\'on anterior (impl\'icitamente) con respecto a $y$, se obtiene que
  $$\frac{\partial z}{\partial y} = \underline{\hspace{5cm}}.$$
  \item[(b)] (12\,\%) Sea $f(x,y)=x^2+k\cdot xy+y^2$. Los valores de $k$ para los cuales la funci\'on $f$ tiene un m\'inimo relativo en $(0,0)$ son:
  $$\underline{\hspace{4cm}}.$$
\end{opciones}

\newpage
\examheader
\vspace{2pt}

\textbf{Preguntas con procedimiento}

\textit{En las preguntas 4 y 5 responda mostrando detalladamente el procedimiento utilizado.}

\bigskip
\noindent\textbf{4.} (16\,\%) Considere la funci\'on
$$
f(x,y)=
\begin{cases}
\dfrac{y^3}{x^2+y^2}+x+y+1, & (x,y)\neq (0,0),\\[6pt]
1, & (x,y)=(0,0).
\end{cases}
$$
Si se sabe que $f_x(0,0)=1$ y $f_y(0,0)=2$.

\begin{opciones}
  \item[(a)] (8\,\%) Determine si $f$ es continua en $(0,0)$.
  \item[(b)] (8\,\%) Determine si $f$ es diferenciable en $(0,0)$.
\end{opciones}

\vfill
\examfooter
\newpage
\examheader
\vspace{2pt}

\noindent\textbf{5.} (14\,\%) Aplique el m\'etodo de Multiplicadores de Lagrange para encontrar los puntos de la hip\'erbola $xy=1$ m\'as cercanos al origen.

\vfill
\examfooter

\end{document}
